%%%%%%%%%%%%%%%%%%%%%%%%%%%%%%%%%%%%%%%%%%%%%%%%%%%%%%%%%%%%%%%%%
\part{Experimentación y resultados}
\label{part:experimentacion}
%%%%%%%%%%%%%%%%%%%%%%%%%%%%%%%%%%%%%%%%%%%%%%%%%%%%%%%%%%%%%%%%%

\section{Experimentación y análisis de resultados}
\label{sec:experimentacion-resultados}

\subsection{Protocolo experimental}

\textit{Versión resumida del protocolo ya definido en la Parte~\ref{part:gestion-tesis}, con referencias cruzadas a las decisiones de la Parte~\ref{part:diseno-sistema}. Se consolidará aquí cuando se cierre el diseño experimental definitivo.}

\subsection{Resultados}

\textit{Capítulo a redactar cuando los experimentos estén cerrados. Se incluirán:} 
\begin{itemize}[nosep,leftmargin=*]
\item Caracterización del teacher (calidad y servicio).
\item Caracterización del student (calidad y servicio).
\item Comparativa de las tres políticas A, B, C en calidad, coste y latencia.
\item Análisis de la frontera de Pareto calidad-coste.
\item Consumo energético medido frente a las estimaciones iniciales de la Parte~\ref{part:gestion-tesis} (sostenibilidad).
\end{itemize}

\subsection{Discusión}

\textit{A redactar con los resultados en la mano: qué se confirma de las hipótesis del diseño, qué sorprende, qué limitaciones salen a la luz.}

%%%%%%%%%%%%%%%%%%%%%%%%%%%%%%%%%%%%%%%%%%%%%%%%%%%%%%%%%%%%%%%%%
